\Needspace{12\baselineskip}
\originalchapter{多孔介质和双扩散}
\originalsection{Darcy定律与均匀化尺度}
孔隙中的流体仍满足局部黏性方程，但宏观观测往往不解析每条孔道。低速、孔隙结构固定时，体积平均速度 $\mathbf q=-(K/\mu)(\nabla p-\rho\mathbf g)$，$K$ 为长度平方量纲的渗透率。$\mathbf q$ 按总体截面积定义，不等于孔内平均流速，二者差一个孔隙率因子。

均匀化需要孔尺度远小于宏观场变化尺度，局部胞问题给渗透率张量；各向同性介质才可用标量 $K$。达西速度的线性关系不是所有高速渗流都有效，惯性较大时可出现非线性修正。薄隙平均的平板流能提供 $K=b^2/12$ 的精确类比，见本卷所选题解。

\originalsection{黏性指进和热对流}
在薄隙或多孔介质中，低黏度流体推进高黏度流体时，略微前伸的界面更易获得流量，会继续前伸。Darcy势流的界面分析给出的增长率有速度、黏度差产生的 $|k|$ 项和表面张力稳定的 $-|k|^3$ 项。稳定性符号取决于推动顺序，不能只见到两种不同黏度就判断必失稳。接触角、薄膜和三维孔隙结构会修正模型。

多孔热层的动量近似即时达西响应，温度仍有时间演化。以层厚 $H$、温差 $\Delta T$ 定义 $Ra_D=g\alpha\Delta T KH/(\nu\kappa)$。对不可穿透、固定温度上下边界，正弦模的中性曲线 $Ra_D=(k^2+\pi^2)^2/k^2$，最小值为 $4\pi^2$、$k=\pi$。这里动量没有自由流体的惯性和黏性拉普拉斯，因此与自由边界Rayleigh--Bénard的 $27\pi^4/4$ 不同。

\originalsection{双扩散的时间竞争}
若温度和盐分同时改变密度，密度扰动为 $\rho'/\rho_0=-\alpha T'+\beta_S S'$。热和盐的扩散系数通常不同，即使总密度梯度稳定，两种场的不同弛豫速度也可能导致增长。在线性模式中，$T'$ 和 $S'$ 分别含 $(\sigma+\kappa_Tk^2)^{-1}$、$(\sigma+\kappa_Sk^2)^{-1}$ 的响应，浮力是二者带不同符号的合成。

盐指和振荡双扩散取决于上下梯度方向、扩散比和流体或多孔动量模型。不能只用总的一个“有效扩散系数”代替两套方程。源课的地热羽流、盐热失稳全文为进一步核对入口，本卷没有求出其完整参数稳定域。

\originalsection{Darcy定律的平均与孔隙速度}
孔内微观流可满足Stokes方程，但复杂孔道不宜逐个直接求解。线性响应意味着宏观单位总面积流量 $\mathbf q=-\mathbf K(\nabla p-\rho\mathbf g)/\mu$，$\mathbf K$ 为渗透率张量、量纲长度平方。各向同性时为标量 $K$。孔内平均速度是 $\mathbf q/\varepsilon_p$，$\varepsilon_p$ 为孔隙率；混用两者会改变示踪旅行时间。

平均需要孔尺度远小于关注尺度、孔内Re小及连通孔隙。靠壁过渡、惯性阻力和多相润湿可超出基本Darcy式。用平行细管模型，平均流量给 $K$ 与孔径平方相关；它说明几何来源而非普适的真实岩石系数。

\originalsection{Hele--Shaw为何等价Darcy}
两平板间距 $b$ 远小于水平尺度，三维压力在厚度上近恒定、速度满足 $\mu u_{zz}=p_x$。无滑移两面得抛物线，厚度平均 $\bar{\mathbf u}=-b^2\nabla_h p/(12\mu)$，等效 $K=b^2/12$。固体随机孔道与确定平间隙的微观结构不同，但在长尺度有同型平均方程。

厚度平均后表面张力、边界湿润和三维端部可能产生修正。热传导如果同时经过玻璃板，使用流体单独热扩散率未必合适。所选原题保留热模型选择，不将一个 $K$ 等价推广成所有热/质输运都完全等价。

\originalsection{Saffman--Taylor指进的线性色散}
水平Darcy/Hele--Shaw中，流体1以均速 $U$ 向右推动流体2，黏度分别 $\mu_1,\mu_2$。压力扰动满足Laplace，界面波数 $k$ 的两侧扰动按 $\e^{-|k||x-Ut|}$ 衰减。运动学界面速度与法向Darcy流匹配，动态条件用张力曲率，消元得增长率
\[
\sigma(k)=U\frac{\mu_2-\mu_1}{\mu_1+\mu_2}|k|
-\frac{\gamma_s K}{\mu_1+\mu_2}|k|^3.
\]
低黏度推动高黏度时第一项为正，小凸起获得更大流速，继续伸出；张力抑制短波。临界波数平方 $k_c^2=U(\mu_2-\mu_1)/(\gamma_sK)$，最大增长 $k_m=k_c/\sqrt3$。反过来黏度顺序全部线性稳定（该水平无重力模型）。实际指形宽度、竞争和尖端选择属非线性，不由这一谱完整决定。

\originalsection{地热羽流的Darcy相似尺度}
二维竖向多孔热羽流，Darcy速度 $w\sim K g\alpha_T\vartheta/\nu$，热方程平衡 $w/z\sim\kappa/\delta^2$，单位宽热通量 $\int w\vartheta\dd x$ 恒定。所以 $\vartheta^2\delta$ 常量、$\delta^2\vartheta\sim z$，合并给 $\vartheta_c,w_c\propto z^{-1/3}$、$\delta\propto z^{2/3}$。与无孔层流羽流的 $w_c\propto z^{1/5}$ 不同，因为动量是阻力平衡而非惯性/黏性随程演化。

该相似形式假定渗透率和热参数常、线热源、单相Boussinesq流。实际地热可能相变、各向异性和热源有限，不能凭这个尺度直接报告某地热井的流量。

\originalsection{Rayleigh--Darcy对流的本征值}
均匀多孔层厚 $H$、下热上冷，取 $Ra_D=g\alpha_T\Delta T K H/(\nu\kappa)$。在本文规定的有效热容量/扩散模型中，线性化并以 $H,H^2/\kappa$ 缩放，温度扰动满足 $\theta_t-w=\nabla^2\theta$，Darcy压力消元给 $\nabla^2w=Ra_D\nabla_h^2\theta$。

不渗透、恒温两面 $w=\theta=0$，取 $\sin(\pi z)\e^{ikx}$、中性 $\sigma=0$，得 $Ra_D=(k^2+\pi^2)^2/k^2$。对 $k^2$ 最小化给 $k_c=\pi,Ra_{D,c}=4\pi^2$。普通无孔RBC自由边界的 $27\pi^4/4$ 来源于黏性四阶动量方程，这里Darcy是二阶，故临界数不同。

\originalsection{双扩散的两个中性机制}
温度底热使密度低、盐度底高使密度高，两者可相互竞争。取热扩散率 $\kappa_T$，盐/热比 $\tau=\kappa_S/\kappa_T$，两Rayleigh参数均以 $\kappa_T$ 定义。单模 $Q=k^2+\pi^2$ 给
\[
(\sigma+Q)\theta=w,\quad(\sigma+\tau Q)s=w,\quad
Qw=k^2(Ra_T\theta-Ra_Ss).
\]
消元是一个二次特征方程
\[
Q\sigma^2+[Q^2(1+\tau)-k^2(Ra_T-Ra_S)]\sigma
+\tau Q^3-k^2Q(\tau Ra_T-Ra_S)=0.
\]
定常中性由常数项零给 $Ra_T-Ra_S/\tau=Q^2/k^2$。振荡中性由一次项零、常数项仍正给 $Ra_T-Ra_S=(1+\tau)Q^2/k^2$。两个阈值不同，扩散速度差允许相位延迟形成振荡；若 $\tau=1$，只剩合成浮力的单扩散问题。

有限层允许离散竖向模，还要比较最先失稳的横向波数。源课全热盐稳定域需要按参数定义实际求这些分支；本节已推导两种机制的判据，但未把未运行的全参数图伪装成数值结果。
